\begin{proof}
The moment identities in
\eqref{eq:p071-support-plane-balanced-data} are equivalent to the ambient
barycentre identity
\[
                  \sum_a\lambda_a(r_av_{i_a},s_aw_{j_a})=0
                  \quad\text{in }E\times F.
\]
Indeed, $v_1,v_2$ are a basis and $v_3=-v_1-v_2$, so vanishing of the first
component says that the three first radial moments are equal.  The same
argument applies to the second component.

We first prove that the five-atom tests imply all finite tests.  Suppose a
balanced family $z_1,\ldots,z_m$ with weights $\lambda_a$ violates
\eqref{eq:p071-support-plane-jensen}, and write
\[
                         x_a=(r_av_{i_a},s_aw_{j_a})\in\mathbf R^4.
\]
Consider the compact polytope
\[
 Q=\left\{\mu\in\mathbf R_{\geq0}^m:
        \sum_a\mu_a=1,\quad \sum_a\mu_ax_a=0\right\}.
\]
It contains $\lambda$.  Minimize
$\sum_a\mu_au(z_a)$ on $Q$ and choose an extreme minimizer $\mu^*$.
The equality matrix defining $Q$ has at most five independent rows.  If
more than five coordinates of $\mu^*$ were positive, the corresponding
columns would be linearly dependent.  A nonzero vector in that dependence,
extended by zero on the remaining coordinates, would give
\[
                         \mu^*+\varepsilon h,\qquad
                         \mu^*-\varepsilon h\in Q
\]
for all sufficiently small $\varepsilon>0$, contradicting extremality.
Thus $\mu^*$ has at most five positive coordinates.  After discarding its
zero coordinates, it is a balanced family of at most five atoms, and
\[
 \sum_a\mu_a^*u(z_a)
 \leq\sum_a\lambda_au(z_a)<u(o).
\]
This contradicts the five-atom tests.  The converse implication is
immediate.

We next pass from the mixed inequalities to one common support plane.  For
$0<\rho<R$, let $K_\rho\subset S\times T$ be the compact product collar
on which both radial coordinates are at most $\rho$.  Its convex hull
contains $0$ as an interior point of $E\times F$: the origin is interior
to each triangle with vertices $\rho v_1,\rho v_2,\rho v_3$ and
$\rho w_1,\rho w_2,\rho w_3$.  On $\operatorname{conv}(K_\rho)$
consider the convex envelope
\[
 f_\rho(x)=
 \inf\left\{
   \sum_a\lambda_au(z_a):
   \lambda_a>0,\ \sum_a\lambda_a=1,\ \sum_a\lambda_az_a=x
 \right\}.
\]
The mixed inequalities and the one-atom measure at $o$ give
$f_\rho(0)=u(o)$.  The function $f_\rho$ is finite and convex, because
$u$ is continuous on the compact set $K_\rho$.  The finite-dimensional
supporting-hyperplane theorem at the interior point $0$ therefore gives
$L_\rho\in(E\times F)^*$ such that
\[
                         f_\rho(x)\geq u(o)+L_\rho(x).
\]
Since $f_\rho(z)\leq u(z)$ for $z\in K_\rho$, the same affine function
minorizes $u$ on $K_\rho$.

It remains to make the linear form independent of the collar.  Put
\[
 \mathcal A_\rho=
 \left\{L\in(E\times F)^*:
       u(z)\geq u(o)+L(z)\text{ for every }z\in K_\rho\right\}.
\]
These sets are nonempty, closed, convex, and decrease as $\rho$ increases.
Fix $\rho_0>0$.  If $L=(A,B)\in\mathcal A_{\rho_0}$, the three first
axis points give
\[
 A(v_i)\leq
 \frac{u(p_i(\rho_0),\eta_T)-u(o)}{\rho_0}
 \qquad(i=1,2,3).
\]
Because $A(v_1)+A(v_2)+A(v_3)=0$, these three upper bounds also give lower
bounds for every $A(v_i)$.  Since $v_1,v_2$ are a basis, $A$ lies in a
bounded subset of $E^*$.  The second-axis points give the same conclusion
for $B$.  Hence $\mathcal A_{\rho_0}$ is compact.

Choose $\rho_n\uparrow R$ with $\rho_n\geq\rho_0$.  The nested nonempty
closed subsets $\mathcal A_{\rho_n}$ of the compact set
$\mathcal A_{\rho_0}$ have nonempty intersection.  Any
$L=(A,B)$ in that intersection gives
\eqref{eq:p071-support-plane-minorant} on the whole open product chart.

Conversely, suppose \eqref{eq:p071-support-plane-minorant} holds and take
balanced data.  Multiplication by the weights and summation give
\[
 \sum_a\lambda_au(z_a)
 \geq u(o)+
 A\left(\sum_a\lambda_ar_av_{i_a}\right)+
 B\left(\sum_a\lambda_as_aw_{j_a}\right)
 =u(o).
\]
Thus the common support plane is equivalent to the mixed double-node
inequalities.

Finally, divide the embedded points in
\eqref{eq:p071-support-plane-sharp-atoms} by $\rho$.  Their coordinates in
$E\times F$ are
\[
 (1,0,0,0),\quad (0,1,0,0),\quad (-1,-1,1,0),\quad
 (0,0,0,1),\quad (0,0,-1,-1).
\]
Their sum is zero.  Taking the last point as origin for affine differences
gives
\[
 \det
 \begin{pmatrix}
 1&0&-1&0\\
 0&1&-1&0\\
 1&1&2&1\\
 1&1&1&2
 \end{pmatrix}
 =5.
\]
Hence the five points are affinely independent.  Their barycentric
coordinates at zero are uniquely $(1/5,\ldots,1/5)$, so no proper
subfamily has barycentre zero.
\end{proof}
